\subsection{Invariantes diferenciales fundamentales de un morfismo de preesquemas}
\label{IV.16.3}

\begin{definition}[16.3.1]
\label{IV.16.3.1}
Sean $f:X\to S$ un morfismo de preesquemas, $\Delta_f:X\to X\times_S X$
el morfismo diagonal correspondiente, que es una inmersión (I, 5.3.9 y
$\mathbf{Err}_{\mathrm{III}}$, 10). Se denota por $\mathscr{P}_f^n$ o
$\mathscr{P}_{X/S}^n$, y se llama haz de partes principales de orden $n$
del $S$-preesquema $X$, al haz de anillos aumentado sobre $\mathcal{O}_X$, $n$-ésimo
invariante normal de $\Delta_f$ (16.1.2). Se pone
$\mathscr{P}_f^\infty=\mathscr{P}_{X/S}^\infty=
\varprojlim_n\mathscr{P}_{X/S}^n$,
$\mathscr{G}r_n(\mathscr{P}_f)=\mathscr{G}r_n(\mathscr{P}_{X/S})=
\mathscr{G}r_n(\Delta_f)$ (16.1.2); el $\mathcal{O}_X$-módulo
$\mathscr{G}r_1(\Delta_f)$, ideal de aumentación de $\mathscr{P}_{X/S}^1$, se
denota también por $\Omega_f^1$ o $\Omega_{X/S}^1$, y se llama el $\mathcal{O}_X$-módulo
de diferenciales de primer orden de $f$, o de $X$ respecto de $S$, o del $S$-preesquema $X$.
\end{definition}

De esta definición resulta que $\mathscr{P}_{X/S}^0$ se identifica
canónicamente con $\mathcal{O}_X$ (16.1.2).

Se tiene (16.1.2.2) un homomorfismo canónico sobreyectivo de
$\mathcal{O}_X$-álgebras graduadas
\[
  \mathscr{S}_{\mathcal{O}_X}^{\bullet}(\Omega_{X/S}^1)
  \longrightarrow\mathscr{G}r_\bullet(\mathscr{P}_{X/S}).
  \tag{16.3.1.1}
  \label{IV.16.3.1.1}
\]
Resulta asimismo de la definición (16.3.1) que, para todo abierto $U$ de $X$, se tiene
$\mathscr{P}_{f|U}^n=\mathscr{P}_f^n|U$,
$\mathscr{P}_{f|U}^\infty=\mathscr{P}_f^\infty|U$,
$\mathscr{G}r_n(\mathscr{P}_{f|U})=\mathscr{G}r_n(\mathscr{P}_f)|U$,
$\Omega_{f|U}^1=\Omega_f^1|U$ (es decir, las nociones introducidas son
locales sobre $X$).

\begin{env}[16.3.2]
\label{IV.16.3.2}
Denotemos por $p_1$, $p_2$ las dos proyecciones canónicas del producto fibrado $X\times_S X$;
como $\Delta_f$ es una sección sobre $X$ de $X\times_S X$ a la vez para los
morfismos $p_1$ y $p_2$, cada uno de estos morfismos define, para todo $n$, un
homomorfismo de haces de anillos $\mathcal{O}_X\to\mathscr{P}_{X/S}^n$,
inverso por la derecha de la aumentación $\mathscr{P}_{X/S}^n\to\mathcal{O}_X$
(16.1.7); también puede decirse que se definen así sobre
$\mathscr{P}_{X/S}^n$ dos estructuras cuasicoherentes de álgebra aumentada sobre
$\mathcal{O}_X$; las estructuras correspondientes de $\mathcal{O}_X$-módulo sobre
$\mathscr{G}r_n(\mathscr{P}_{X/S})$ son las mismas. Se tienen asimismo, al pasar al
límite, dos estructuras de $\mathcal{O}_X$-álgebra sobre
$\mathscr{P}_{X/S}^\infty$.
\end{env}

\begin{env}[16.3.3]
\label{IV.16.3.3}
El morfismo $s=(p_2,p_1)_S:X\times_S X\to X\times_S X$ es un automorfismo
involutivo de $X\times_S X$, llamado la simetría canónica, tal que
\[
  p_1\circ s=p_2,\qquad p_2\circ s=p_1,\qquad s\circ\Delta_f=\Delta_f.
  \tag{16.3.3.1}
  \label{IV.16.3.3.1}
\]
Si se pone $s=(\rho,\lambda)$, $p_i=(\pi_i,\mu_i)$ ($i=1,2$),
$\Delta_f=(\delta,\nu)$, $\lambda^\#$ es por tanto un isomorfismo de
$\rho^*(\pi_1^*(\mathcal{O}_X))$ en $\pi_2^*(\mathcal{O}_X)$, y
$\delta^*(\lambda^\#)$ deja invariantes $\delta^*(\mathcal{O}_{X\times_S X})$
y el núcleo $\mathscr{I}$ del homomorfismo
$\nu^\#:\delta^*(\mathcal{O}_{X\times_S X})\to\mathcal{O}_X$. Por consiguiente:
\end{env}

\begin{proposition}[16.3.4]
\label{IV.16.3.4}
El homomorfismo $\sigma=\delta^*(\lambda^\#)$ deducido de $s$ (y llamado también
simetría canónica) es un automorfismo involutivo del sistema proyectivo
$(\mathscr{P}_{X/S}^n)$ de haces
\oldpage[IV-4]{15}
de anillos aumentados sobre $\mathcal{O}_X$, y por consiguiente también de su límite proyectivo
$\mathscr{P}_{X/S}^\infty$. Este automorfismo intercambia las dos estructuras de
$\mathcal{O}_X$-álgebra sobre los $\mathscr{P}_{X/S}^n$ y sobre
$\mathscr{P}_{X/S}^\infty$.
\end{proposition}

\begin{env}[16.3.5]
\label{IV.16.3.5}
En lo que sigue, las dos estructuras de $\mathcal{O}_X$-álgebra definidas sobre los
$\mathscr{P}_{X/S}^n$ y sobre $\mathscr{P}_{X/S}^\infty$ desempeñarán papeles muy
distintos: convendremos en adelante, salvo mención expresa de lo contrario,
que cuando $\mathscr{P}_{X/S}^n$ o $\mathscr{P}_{X/S}^\infty$ se considere
como una $\mathcal{O}_X$-álgebra, se tratará de la estructura de álgebra definida por
$p_1$.
\end{env}

Para todo abierto $U$ de $X$ y toda sección $t\in\Gamma(U,\mathcal{O}_X)$, se
denotará simplemente por $t.1$ o incluso por $t$ la imagen de $t$ mediante el homomorfismo
estructural $\Gamma(U,\mathcal{O}_X)\to\Gamma(U,\mathscr{P}_{X/S}^n)$ (resp.
$\Gamma(U,\mathcal{O}_X)\to\Gamma(U,\mathscr{P}_{X/S}^\infty)$)
(es decir, el homomorfismo correspondiente a $p_1$).

\begin{definition}[16.3.6]
\label{IV.16.3.6}
Se denota por $d_f^n$, o $d_{X/S}^n$ (resp. $d_f^\infty$, o
$d_{X/S}^\infty$), o simplemente $d^n$ (resp. $d^\infty$), el homomorfismo de
haces de anillos $\mathcal{O}_X\to\mathscr{P}_f^n=\mathscr{P}_{X/S}^n$
(resp. $\mathcal{O}_X\to\mathscr{P}_f^\infty=\mathscr{P}_{X/S}^\infty$)
deducido de $p_2$. Para todo abierto $U$ de $X$, y todo
$t\in\Gamma(U,\mathcal{O}_X)$, $d^nt$ (resp. $d^\infty t$) se llama la parte
principal de orden $n$ (resp. parte principal de orden infinito) de $t$. Se pone
$dt=d^1t-t$, y se dice que $dt$ es el diferencial de $t$ (elemento de
$\Gamma(U,\Omega_{X/S}^1)$, denotado también por $d_{X/S}(t)$).
\end{definition}

De esta definición se sigue inmediatamente que
\[
  d(t_1t_2)=t_1dt_2+t_2dt_1
  \tag{16.3.6.1}
  \label{IV.16.3.6.1}
\]
para cualesquiera $t_1,t_2$ en $\Gamma(U,\mathcal{O}_X)$; es decir, $d$
es una derivación del anillo $\Gamma(U,\mathcal{O}_X)$ con valores en el
$\Gamma(U,\mathcal{O}_X)$-módulo $\Gamma(U,\Omega_{X/S}^1)$.

En todas las notaciones introducidas en (16.3.1) y (16.3.6), se reemplazará
a veces $S$ por $A$ cuando $S=\operatorname{Spec}(A)$.

\begin{env}[16.3.7]
\label{IV.16.3.7}
Supongamos, en particular, que $S=\operatorname{Spec}(A)$ y
$X=\operatorname{Spec}(B)$ son esquemas afines, siendo entonces $B$ una
$A$-álgebra. Entonces $\Delta_f$ corresponde al homomorfismo canónico sobreyectivo
$\pi:B\otimes_A B\to B$ tal que $\pi(b\otimes b')=bb'$, cuyo núcleo es
$\mathfrak{I}=\mathfrak{I}_{B/A}$ (0, 20.4.1); $\mathscr{P}_f^n$ es el haz
estructural del preesquema $\operatorname{Spec}(P_{B/A}^n)$, donde
\[
  P_{B/A}^n=(B\otimes_A B)/\mathfrak{I}^{n+1};
\]
$\mathscr{G}r_\bullet(\mathscr{P}_f)$ es el $\mathcal{O}_X$-módulo
cuasicoherente correspondiente al $B$-módulo graduado
\[
  \operatorname{gr}_{\mathfrak{I}}^\bullet(B\otimes_A B)
  =\bigoplus_{n\geq 0}(\mathfrak{I}^n/\mathfrak{I}^{n+1});
\]
en particular, $\Omega_f^1=\Omega_{X/S}^1$ es el $\mathcal{O}_X$-módulo
cuasicoherente correspondiente al $B$-módulo de diferenciales de primer orden de $B$ con
respecto a $A$, $\Omega_{B/A}^1$ (0, 20.4.3). Los morfismos de proyección
$p_1:X\times_S X\to X$, $p_2:X\times_S X\to X$ corresponden a los dos
homomorfismos de anillos $j_1:B\to B\otimes_A B$, $j_2:B\to B\otimes_A B$
tales que $j_1(b)=b\otimes 1$, $j_2(b)=1\otimes b$, de modo que (por la
convención de (16.3.5)), $P_{B/A}^n$ se considera siempre como una
$B$-álgebra mediante el homomorfismo compuesto
$B\xrightarrow{j_1}B\otimes_A B\to P_{B/A}^n$; el homomorfismo de anillos
$B\xrightarrow{j_2}B\otimes_A B\to P_{B/A}^n$ se denota por $d_{B/A}^n$ y corresponde
a $d_{X/S}^n$, que actúa sobre $\Gamma(X,\mathcal{O}_X)$; para todo $t\in B$, $dt$
es igual a $d_{B/A}t$, definido en (0, 20.4.6) (cf. $\mathbf{Err}_{\mathrm{IV}}$,
11).

Si $\pi_n:B\otimes_A B\to P_{B/A}^n$ es el homomorfismo canónico, se tiene, por tanto,
en virtud de las definiciones precedentes,
\[
  \pi_n(b\otimes b')=b.\pi_n(1\otimes b')=b.d_{B/A}^n(b')
  \quad\text{para }b\in B,\ b'\in B.
  \tag{16.3.7.1}
  \label{IV.16.3.7.1}
\]
\end{env}

\oldpage[IV-4]{16}
\begin{proposition}[16.3.8]
\label{IV.16.3.8}
La imagen del homomorfismo $d_{X/S}^n:\mathcal{O}_X\to\mathscr{P}_{X/S}^n$
genera el $\mathcal{O}_X$-módulo $\mathscr{P}_{X/S}^n$.
\end{proposition}

Se reduce inmediatamente al caso en que $X=\operatorname{Spec}(B)$ y
$S=\operatorname{Spec}(A)$ son afines, y la proposición resulta de (16.3.7.1),
puesto que $\pi_n$ es sobreyectivo. Obsérvese que, en general, $d_{X/S}^n$ no es
sobreyectivo (incluso para $n=1$).

\begin{proposition}[16.3.9]
\label{IV.16.3.9}
Supongamos que $f:X\to S$ es un morfismo localmente de tipo finito. Entonces los
$\mathscr{P}_f^n$ y los $\mathscr{G}r_n(\mathscr{P}_f)$ son
$\mathcal{O}_X$-módulos cuasicoherentes de tipo finito.
\end{proposition}

Esto resulta de (16.1.6) y del hecho de que $\Delta_f$ es localmente de presentación
finita (I, 4.3.1).
